\printedpage{109}
% head: A. Grothendieck : Dérivateurs, Ch. XVIII, Édition des Universités de Montpellier II et Paris VII

On the other hand, each $\underline{S}_\alpha$ being stable in $\mathcal{M}$ under $\varinjlim$ large relative to $\pi'_0$, the same holds for their intersection $\underline{S}$. A fortiori, the latter is stable under $\varinjlim$ large relative to $\pi$, so the canonical fully faithful functor $\Ind_\pi(\underline{S}(\pi))\longrightarrow\mathcal{M}$ factors through $S$ as
\[ \Ind_\pi\underline{S}(\pi) \incl \underline{S} \; , \]
we must prove that this functor is essentially surjective. So let $x\in S=\Ob\underline{S}$, it all comes down to proving that
\[ \varinjlim_{i\in\underline{S}(\pi)/x} x_i \xrightarrow{\ \sim\ } x \; , \tag*{$(*)$} \]
and that $\underline{S}(\pi)/x$ is large relative to $\pi$.

Note that, since $\pi'_0\geqslant\pi_0$, and $\pi^{\pi'_0}=\pi$, we will have $\pi^{\pi_0}=\pi$, so $\pi$ is a useful cardinal for $\mathcal{M}$, so $I=\mathcal{M}_\pi/x$ is large relative to $x$, and we have % [sic] grande devant $x$ -> grande devant $\pi$, apparently intended (one speaks of being large before a cardinal)
\[ \varinjlim_{i\in I=\mathcal{M}_\pi/x} x_i \xrightarrow{\ \sim\ } x \; . \tag*{$(**)$} \]
Thus, if we admit that $\underline{S}(\pi)/x$ is large relative to $\pi$,
\msp{page 294}
the fact that $(*)$ is an isomorphism is equivalent to saying that the canonical functor
\[ J \overset{\text{def}}{=} \underline{S}(\pi)/x \incl \mathcal{M}_\pi \; = \; I \] % [sic] $\mathcal{M}_\pi = I$ -> $\mathcal{M}_\pi/x = I$, since $I=\mathcal{M}_\pi/x$ above and in the display for $J_\alpha$ below
is \textit{cofinal} (cf. 4.4.2, p. 79). But we see at once that this functor is an inclusion functor of a strictly full subcategory (stable under $\varinjlim$ large relative to $\pi'_0$, of cardinality $\leqslant\pi$).

Note now that the arrow $(**)$ similar to $(*)$, with $\underline{S}(\pi)$ replaced by $\underline{S}_\alpha(\pi)$, is indeed an isomorphism, since $x\in\Ob\underline{S}\subset\Ob\underline{S}_\alpha$, and $\underline{S}_\alpha$ is $\pi$-accessible. This is interpreted as the fact that the inclusion functor
\[ J_\alpha \overset{\text{def}}{=} \underline{S}_\alpha(\pi)/x \incl \mathcal{M}_\pi/x \; = \; I \]
is cofinal. This is again an inclusion of a strictly full subcategory $J_\alpha$ of $I$, stable under $\varinjlim$ large relative to $\pi'_0$ and of cardinality $\leqslant\pi$, and we have
\[ J \; = \; \bigcap J_\alpha \; . \]

Thus the desired condition will follow from the
\msp{page 295}

% marginal hand-drawn mark (manuscript) in the left margin beside the following lemma
\textbf{Lemma 4.13.10} (\footnote[125]{We must assume that $\pi$ is not a limit cardinal, to be able to conclude that $I$ is stable under $\varinjlim$ of type $I_{\alpha_\pi}$, and the $J_r$ stable in $I$ under these limits.},\footnote[126]{Cf. more general form of the lemma in 4.13.10.2, p. 302.}). \textit{Let $\pi'_0$ and $\pi>\pi'_0$ be infinite cardinals, $I$ a category equivalent to a small category, $(J_\alpha)_{\alpha\in A}$ a family of subcategories $J_\alpha$ of $I$. We assume:}
\begin{itemize}
\item[\cercle{a}] \textit{$I$ is large relative to $\pi$ (}\footnote[127]{It suffices that $I$ be large relative to every $\pi'<\pi$.}\textit{).} % footnote 127 printed at the foot of p. 110
\item[\cercle{b}] \textit{$I$ is stable under $\varinjlim$ of cardinality $\leqslant\pi$, and large relative to $\pi'_0$ (}\footnote[128]{Here stability under $\varinjlim$ of type $I_{\alpha_\pi}$ suffices, cf. page 302 [?].}\textit{).} % footnote 128 printed at the foot of p. 110
\end{itemize}
% foot: Version 17 décembre 2025

\printedpage{110}
% head: A. Grothendieck : Dérivateurs, Ch. XVIII, Édition des Universités de Montpellier II et Paris VII
\begin{itemize}
\item[\cercle{c}] \textit{$\forall\ r\in A$, $J_r$ is stable in $I$ under the direct limits envisaged in b) (in particular, $J_r$ is a} strictly full \textit{subcategory of $I$).} % [sic] envisagés -> envisagées, agreement with « limites »
\item[\cercle{d}] \textit{$\forall\, r\in A$, $J_r$ is a} cofinal \textit{subcategory of $I$ (}large relative to $\pi$\textit{).}
\item[\cercle{e}] $\card A\leqslant\pi$.
\end{itemize}
\textit{Let $J=\bigcap_{r\in A}J_r$. Then $J$ is still cofinal in $I$, and large relative to $\pi$ (and moreover stable in $I$ under $\varinjlim$ of cardinality $\leqslant\pi$ and large relative to $\pi'_0$).}

In other words, if $I$ is large relative to $\pi$, the intersection $J$ of a family $(J_\alpha)_{\alpha\in A}$ of strictly full subcategories cofinal and large relative to $\pi$, is still
\msp{page 296}
cofinal and large relative to $\pi$, \textit{provided that} $\card A\leqslant\pi$, and that we postulate a stability property of $I$ and of the $J_\alpha$ in $I$, for direct limits of cardinality $\leqslant\pi$ and large relative to a suitable given cardinal $\pi'_0<\pi$.

\textbf{4.13.10.1.} Note moreover that if $f : J\longrightarrow I$ is a cofinal \textit{fully faithful} functor, with $I$ large relative to $\pi$, then $J$ is automatically large relative to $\pi$. This follows e.g. from 4.9.6.3 2° (p. 225) by the criterion $\mathrm{c}_\pi$ (since for $f$ fully faithful, condition F $3_\pi$ is tautologically satisfied). We can also see it thus: Let $B$ be a set with $\card B\leqslant\pi$, we must prove that $J\longrightarrow J^B$ is cofinal, now we have
\[
\begin{tikzcd}
J \arrow[r, "\text{cof.\,?}"] \arrow[d, "\text{cof.}"'] & J^B \arrow[d, "\text{cof.}"] \\
I \arrow[r, "\text{cof.}"] & I^B \; ,
\end{tikzcd}
\]
where three of the four arrows of the square are cofinal, $J^B\longrightarrow I^B$ being so as a product of cofinal functors between \textit{filtering} categories (cf. 4.9.6.5, page 227) -- indeed we already know that $J$ is filtering.

\msp{page 297}
Thus the composite
\[ J \longrightarrow J^B \xrightarrow{\text{cof.}} I^B \]
is cofinal. But since the second functor $J^B\longrightarrow I^B$ is fully faithful, it follows at once that $J\longrightarrow J^B$ is cofinal (for $\xi\backslash J$ for $\xi\in\Ob J^B$ is the same, up to isomorphism, whether we consider $\xi$ as an object of $J^B$, or of $I^B$).

Thus in the statement of condition d) of the lemma, it is unnecessary to specify that $J_\alpha$ is large relative to $\pi$, and likewise in the conclusion about $J$.

Since $J\incl I$ is fully faithful, to see that this functor is cofinal, it suffices to see that it satisfies F 1, i.e. that for every $i\in\Ob I$, there exists $j\in J$ with $i\longrightarrow j$. We are going to construct $j$ as the limit of an ordinal direct system of objects $i_\alpha$ of $I$, indexed by the ordinals $\alpha$ such that $\card\alpha<\pi$. Denote by $\alpha_\pi$ the smallest ordinal such that % [sic] le plus petite ordinal -> le plus petit ordinal, agreement with « ordinal »
% foot: Version 17 décembre 2025
\printedpage{111}
% head: A. Grothendieck : Dérivateurs, Ch. XVIII, Édition des Universités de Montpellier II et Paris VII
\msp{page 298}
% marginal hand-drawn mark (manuscript) in the left margin beside the following two lines
$\card\alpha_\pi=\pi$. The ordered set $I_{\alpha_\pi}$ of ordinals $\alpha<\alpha_\pi$ being large relative to $\pi'_0$ (\footnote[129]{false if $\pi$ is a limit cardinal.}) and of cardinality $\pi$, we will therefore be able to take the limit $\varinjlim\limits_{\alpha\in I_{\alpha_0}} i_\alpha$ in $I$. % [sic] $I_{\alpha_0}$ -> $I_{\alpha_\pi}$, the index set just defined is $I_{\alpha_\pi}$

Construction of $(i_\alpha)_{\alpha\in I_{\alpha_0}}$. Let % [sic] $I_{\alpha_0}$ -> $I_{\alpha_\pi}$, the index set just defined is $I_{\alpha_\pi}$
\[ \alpha_\pi \; = \; \text{smallest ordinal } \alpha \text{ such that } \card\alpha=\pi, \tag*{(1)} \]
so this is the ordinal characterized by
\[ \left\{
\begin{aligned}
&\card\alpha_\pi \; = \; \pi \\
&\alpha<\alpha_\pi \;\; \Longrightarrow \;\; \card\alpha \; < \; \pi \; ,
\end{aligned}
\right. \tag*{(2)} \]
% marginal hand-drawn mark (manuscript) in the left margin beside the following paragraph
or again by the fact that the ordered set $I_{\alpha_\pi}$ of ordinals $\alpha<\alpha_\pi$ is \textit{perfectly well ordered}. Thus, $I_{\alpha_\pi}$ \textit{is large relative to $\pi'$ for every cardinal $\pi'<\pi$} (\footnote[130]{False if $\pi$ is a limit cardinal.}). In conditions b) c) of 4.3.10, it will suffice to assume $I$ stable under $\varinjlim$ of type $I_{\alpha_\pi}$, and the $J_r$ ($r\in A$) stable in $I$ under this type of limits. % [sic] 4.3.10 -> 4.13.10, the conditions b) c) are those of Lemme 4.13.10

\textbf{Lemma 4.13.10.1.} \textit{Let $\pi$ be an infinite cardinal, $A$ a nonempty set such that} $\card A\leqslant\pi$. \textit{Then}
\msp{page 299}
\textit{there exists a map}
\[ \varphi \; : \; I_{\alpha_\pi} \longrightarrow A \]
\textit{such that for every $r\in A$, $\varphi^{-1}(\{r\})$ is a cofinal subset of $I_{\alpha_\pi}$, or what is more precise,} $\card\varphi^{-1}(\{r\})=\pi$. % [sic] ce que est -> ce qui est, relative pronoun subject

Indeed, it suffices to consider a set $A'$ of cardinality $\pi$ (e.g. $A'=I_{\alpha_\pi}$), the map
\[ \psi \; = \; \mathrm{pr}_1 \; : \; A\times A'\longrightarrow A \; , \]
and to give oneself a bijection
\[ A\times A' \xrightarrow{\ \sim\ } I_{\alpha_\pi} \]
which exists since
\[ \card\,(A\times A') \; = \; \card A\cdot\card A' \; = \; \card A' \; = \; \pi \; = \; \card I_{\alpha_\pi} \; , \]
since $0\neq\card A\leqslant\card A'=\pi$ and $\pi$ infinite. We assume $\varphi$ chosen as in 4.13.10.1, and we will construct the direct system $(i_\alpha)_{\alpha\in I_{\alpha_\pi}}$ in such a way that we have
\[ i_\alpha \; \in \; J_{\varphi(\alpha)} \; . \]

Setting
\[ j \; = \; \varinjlim_{\alpha\in I_{\alpha_\pi}} i_\alpha \; , \]
% foot: Version 17 décembre 2025
\printedpage{112}
% head: A. Grothendieck : Dérivateurs, Ch. XVIII, Édition des Universités de Montpellier II et Paris VII
\msp{page 300}
for every $r$, since
\[ I_r \overset{\text{def}}{=} \varphi^{-1}(r) \; \subset \; I_{\alpha_\pi} \]
is cofinal in $I_{\alpha_\pi}$, we will therefore have
\[ j \; = \; \varinjlim_{\alpha\in I_r} i_\alpha \; , \]
and by virtue of $(*)$, the $i_\alpha$ for $\alpha\in I_r$ are in $J_r$, so by passage to the limit, using hypothesis c), we find that $j\in J_r$, so
\[ j \; \in \; \bigcap_{r\in A} J_r \; = \; J \; , \]
and since we will choose
\[ i_0 \; = \; i \; , \tag*{$(**)$} \]
we will indeed have found $i\longrightarrow j$ with $i\in J$. So it suffices to construct the direct system $(i_\alpha)$ satisfying $(*)$ and $(**)$. % [sic] avec $i\in J$ -> avec $j\in J$, it is $j$ that was shown to lie in $J$

If the direct system has been constructed up to the stage $i_\alpha\in\Ob J_{\varphi(\alpha)}$, we continue with
\[ i_\alpha \longrightarrow i_{\alpha+1} \; \in \; \Ob J_{\varphi(\alpha+1)} \; , \]
which is possible since $J_{\varphi(\alpha+1)}$ is
\msp{page 301}
cofinal. So it remains to construct $i_\alpha$ for $\alpha$ a limit ordinal. Now the set of $i_\beta$ for $\beta<\alpha$ is of cardinality $<\pi$, and since by a) $I$ is large relative to $\pi$ (it suffices that it be so relative to every cardinal $<\pi$), there exists an $i\in\Ob I$ which dominates all the $i_\beta$. Let us choose
\[ i_\beta \xrightarrow{v_\beta} i \; . \]

We will not have a priori commutativity in
\[
\begin{tikzcd}
i_{\beta'} \arrow[r, "\tau_{\beta,\beta'}"] \arrow[rr, bend right=15, "v_{\beta'}"'] & i_\beta \arrow[r, "v_\beta"] & i \; ,
\end{tikzcd}
\] % doubt resolved against PDF; see corrections.csv
but since the set of double arrows $i_{\beta'}\mathrel{\overset{v_\beta\tau_{\beta,\beta'}}{\underset{v_{\beta'}}{\rightrightarrows}}} i$ is of cardinality $<\pi$, $I$ being large relative to $\pi$, there exists $i\longrightarrow i'$ which equalizes them all. Since $J_{\varphi(\alpha)}$ is cofinal in $I$, there exists $i''\in J_{\varphi(\alpha)}$, and $i\longrightarrow i''$. We take $i_\alpha=i''$, with the composite homomorphisms
\[
\begin{tikzcd}
i_\beta \arrow[r, "v_\beta"] \arrow[rrr, dotted, bend right=12, "\tau_{\alpha,\beta}"'] & i \arrow[r] & i' \arrow[r] & i''
\end{tikzcd}
\] % doubt resolved against PDF; see corrections.csv
% foot: Version 17 décembre 2025
\printedpage{113}
% head: A. Grothendieck : Dérivateurs, Ch. XVIII, Édition des Universités de Montpellier II et Paris VII
\msp{page 302}
as transition morphisms. This completes the proof of 4.13.10 and thereby of 4.13.9, 4.13.8.

\textbf{Remark 4.13.10.2.} Here is what we really had to use, to establish the conclusion of 4.13.10, in terms of $I$, $\pi$, $(J_r)_{r\in A}$ (and without introducing any $\pi'_0$).
\begin{itemize}
\item[\cercle{a$'$}] $I$ is large relative to every cardinal $\pi'<\pi$.
\item[\cercle{b$'$}] $I$ is stable under $\varinjlim$ of type $I_{\alpha_\pi}$.
\item[\cercle{c$'$}] $\forall\, r\in A$, $J_r$ is stable in $I$ under $\varinjlim$ of type $I_{\alpha_\pi}$.
\item[\cercle{d$'$}] $\forall\, r\in A$, $J_r$ is cofinal in $I$.
\item[\cercle{e$'$}] $\card A\leqslant\pi$.
\end{itemize}
(\textbf{N.B.} d$'$) and e$'$) are none other than d) and e) of 4.3.10, whereas a$'$), b$'$), c$'$) are weakened versions of a), b), c) respectively.) % [sic] 4.3.10 -> 4.13.10, the conditions d) e) are those of Lemme 4.13.10

\textbf{4.13.8 (reformulation).} \textit{Let, for $r\in A$, $\pi''_r\geqslant\pi'_r\geqslant\pi_0$ (= given useful cardinal for $\mathcal{M}$) be as specified in 4.13.2 b) (and c), let $\pi'=\sup_{r\in A}\pi'_r$, $\pi''=\sup_{r\in A}\pi''_r$, so that for every cardinal $\pi$ with $\pi^{\pi'_0}=\pi\geqslant\pi''$, $\underline{S}_r$ is $\pi$-accessible in $\mathcal{M}$ (and $\pi$ useful for $\underline{S}_r$} and \textit{for $\mathcal{M}$). Let $c$ be a} non-limit \textit{cardinal such that $c\geqslant$} $\card A$ \textit{and $c>\pi'$ (e.g. the successor of} $\sup(\card A,\pi')$\textit{). Let us take $\pi$ such that}
\[ \pi^{\pi'} \; = \; \pi \; \geqslant \; \sup(\pi'',c) \; . \]
\textit{Then $\underline{S}=\bigcap_{r\in A}\underline{S}_r$ is $\pi$-accessible.}

Indeed, if $\alpha_c$ is the first ordinal of cardinality $c$, since $c$ is not a limit, $I_{\alpha_c}$ is large relative to every cardinal $c'<c$, in particular relative to $\pi'$. Now for every $x\in\Ob\mathcal{M}$, $I=\mathcal{M}_\pi/x$ and the $J_r=\underline{S}_r/x$ are stable under $\varinjlim$ large relative to $\pi'$ and of cardinality $\leqslant\pi$,
\msp{page 302 bis}
and the $J_r$ stable in $I$ for these limits, so this is in particular the case for $\varinjlim$ of type $I_{\alpha_c}$. We can therefore apply result 4.13.10.2 to conclude that if the $J_r$ are cofinal in $I$ (which is large relative to $\pi$, $\pi$ being useful for $\mathcal{M}$), then $J=\underline{S}/x$ is so as well. This then implies, as we have seen, that $\underline{S}$ is $\pi$-accessible in $\mathcal{M}$.

\textbf{4.13.10.3.} We can reformulate the preceding reformulation, using the terminology introduced below (4.13.23):

\textit{Let $c$ be a cardinal, $\alpha_c$ the smallest ordinal $\alpha$ such that} $\card\alpha=c$, \textit{$\pi$ a cardinal $\geqslant c$. We assume}
\begin{itemize}
\item[$(*)$] \textit{$\pi$ is useful for $\mathcal{M}$, and $\mathcal{M}$ is stable under $\varinjlim$ of type $I_{\alpha_c}$. (This second condition is the case e.g. if $c$ is} non-limit \textit{and $c>\pi_0$, $\pi_0$ being a cardinal such that $\mathcal{M}$ satisfies \textit{2}$_{\pi_0}$, e.g. a} useful \textit{cardinal for $\mathcal{M}$ (}\footnote[131]{For then $I_{\alpha_c}$ is large relative to $\pi_0$.}\textit{)).}
\end{itemize}
\textit{Then the set of subcategories $\underline{S}$ of $\mathcal{M}$ which are $\alpha_c$-$\pi$-accessible in $\mathcal{M}$ (i.e. which are $\pi$-accessible and stable in $\mathcal{M}$ under $\varinjlim$ of type $I_{\alpha_c}$) is stable under intersection of cardinality $\leqslant c$. More generally, if $\Phi$ is a set of categories such that} % [sic] accessible -> accessibles, agreement with « sous-catégories »
\begin{itemize}
\item[\textit{1°)}] \textit{$I\in\Phi\Longrightarrow$} $\card I\leqslant\pi$ \textit{and}
\item[\textit{2°)}] \textit{$I_{\alpha_c}\in\Phi$,}
\end{itemize}
\textit{then the set of subcategories of $\mathcal{M}$ which are $\Phi$-$\pi$-accessible in $\mathcal{M}$, is stable under intersection of cardinality $\leqslant c$.}
% foot: Version 17 décembre 2025